\section{Discussion}
\label{sec:discussion}

\subsection{Mechanisms of Multimodal Interference and Feature Dilution}
The empirical results across the BraTS2021 modality scaling matrix ($M \in \{1, 2, 3, 4\}$) offer critical mechanistic insights into multimodal representation learning. 
At $M=3$ (T1, T2, FLAIR), all capacity-matched models achieve stable segmentation performance (0.6977--0.7078 DSC). 
However, when the fourth modality ($t_{1\text{ce}}$, gadolinium contrast enhancement) is introduced at $M=4$, `BraTSClassicalBaseline` undergoes a severe performance drop to 0.6536 DSC ($-4.41\%$ loss, with early stopping triggered at Epoch 4). 
Under tight parameter budgets, direct classical pooling across four distinct physical intensity distributions induces feature saturation and gradient interference: the high-intensity enhancing tumour core in $t_{1\text{ce}}$ conflicts with the fluid-attenuated edema representations in FLAIR and tissue contrast in T1/T2. 
In contrast, `BraTSAdaptiveCQ` completely resists this collapse (0.7015 DSC, +4.79\% advantage over classical baseline). 
The gating router dynamically rebalances representation allocation to $\bar{\alpha} = 0.4883$, demonstrating that hybrid adaptive allocation acts as an architectural firewall against cross-modal feature dilution.

\subsection{FiLM Modulation and NISQ Noise Absorption}
Our experiments under open-system gate depolarizing noise ($p \in [0.000, 0.050]$) clarify how classical-quantum hybrid networks behave under hardware degradation. 
In `FixedCQ`, because 100\% of the latent bottleneck flows through the quantum circuit, increasing depolarizing noise attenuates Pauli-Z expectation values towards zero ($\langle Z_i \rangle \to 0$), resulting in a monotonic loss of $-0.32\%$ DSC. 
In `AdaptiveCQ`, the parallel classical surrogate branch $z_C$ acts as an uncorrupted physical buffer, restricting zero-shot degradation to $-0.01\%$ DSC across the entire noise regime. 

Furthermore, our analysis explains why `FixedCQ` does not suffer total failure under depolarizing noise:
In the spatial modulation layer $F' = \gamma(z_Q) \odot F + \beta(z_Q)$, the linear projection heads ($\text{Linear}_{8 \to 32}$) possess learnable bias vectors $\beta_0, \gamma_0$. 
When depolarizing noise contracts $\langle Z_i \rangle \to 0$, the network absorbs this signal attenuation into its linear bias parameters, allowing the downstream 3D CNN decoder (17,521 parameters) to preserve baseline spatial reconstruction. 
This demonstrates that in FiLM-modulated hybrid networks, the classical decoder provides inherent structural stabilization against quantum hardware errors.

\subsection{Resolving the Three Architectural Pitfalls}
This study resolves the three core failure modes that historically undermined claims in hybrid medical QML:
\begin{enumerate}
    \item \textbf{Classical Collapse:} In unregularized hybrid networks, classical layers provide smoother gradients with larger magnitudes, causing routing gates to settle permanently at $\alpha = 0$. By introducing piecewise linear annealed regularization ($\lambda(t)(\bar{\alpha}-0.5)^2$), Q-FUSION enforces bidirectional exploration during early epochs, stabilizing $\bar{\alpha}$ between $0.47$ and $0.52$.
    \item \textbf{Spatial Reconstruction Paradox:} Bypassing high-capacity 200,000-parameter unflattening layers, our 576-parameter FiLM module proves that compact 8-dimensional quantum states can effectively modulate 3D spatial feature maps without sacrificing boundary fidelity.
    \item \textbf{Classical Control Parity:} By evaluating an exact 24-parameter capacity-matched classical surrogate (`AdaptiveCC`, 0.0000\% capacity variance), we ensure that reported gains are not artifacts of parameter inflation.
\end{enumerate}

\subsection{Non-Hype Boundaries and Limitations}
In accordance with ethical scientific reporting, we explicitly delimit the scope of our findings:
First, we do not claim unconditional quantum supremacy or exponential speedups over unconstrained classical models. Our contribution demonstrates that \textit{under an identical, constrained parameter budget, adaptive classical-quantum allocation outperforms rigid quantum designs and mitigates cross-modal interference}.
Second, we do not claim that Q-FUSION creates fault-tolerant quantum hardware or solves physical decoherence. Rather, we demonstrate that adaptive routing provides architectural resilience against NISQ noise by dynamically adjusting representation weights.
Finally, all quantum simulations were evaluated on statevector and density-matrix emulators. Validation on physical superconducting transmon or trapped-ion hardware across multi-patient cohorts remains an important direction for future investigation.
